En conclusión, los resultados experimentales muestran que la notación asintótica permite anticipar correctamente la tendencia general de los algoritmos, 
pero no explica por sí sola su comportamiento real. En ordenamiento, \texttt{Mergesort} y \texttt{Quicksort} presentaron un rendimiento estable y cercano a lo esperado, 
mientras que \texttt{std::sort} evidenció una mayor sensibilidad a la implementación y a la medición de la memoria. En multiplicación de matrices, \texttt{Naive} resultó 
más conveniente para tamaños pequeños y medianos por su simplicidad y bajo consumo de memoria, mientras que \texttt{Strassen} mostró el costo asociado a su implementación recursiva, 
especialmente en uso de memoria. En conjunto, este trabajo permite concluir que la eficiencia práctica de un algoritmo no depende únicamente de su complejidad teórica, sino 
también de factores como la implementación, el tamaño real de la entrada y, adentrando mucho más en este análisis, el método de medición de memoria que pudo ser evidenciado como un factor
gatillante a obtener resultados que difieren a lo teórico.
